\documentclass[presentation,10pt,aspectratio=169]{beamer}

% --- Essential Packages ---
\usepackage[utf8]{inputenc}
\usepackage{bm}
\usepackage{amsmath}
\usepackage{amssymb}
\usepackage{booktabs}
\usepackage{tikz}

% --- Theme and Visual Setup ---
\usetheme{Madrid}
\usecolortheme{whale}
\useinnertheme{rounded}
\useoutertheme{infolines}
\usefonttheme{professionalfonts}

\setbeamertemplate{navigation symbols}{}

% --- Professional Color Palette ---
\definecolor{UdeAGreen}{RGB}{0, 111, 61}
\definecolor{OdysseyBlue}{RGB}{28, 117, 188}
\definecolor{PhysicsOrange}{RGB}{230, 108, 44}
\definecolor{StructureGray}{RGB}{80, 80, 80}

\setbeamercolor{titlelike}{parent=palette primary,fg=white,bg=UdeAGreen}
\setbeamercolor{block title}{bg=UdeAGreen,fg=white}
\setbeamercolor{block body}{bg=black!5,fg=black}
\setbeamercolor{block title example}{bg=OdysseyBlue,fg=white}
\setbeamercolor{block body example}{bg=black!5,fg=black}

% --- Metadata Preamble ---
\title[Hybrid Symplectic Propagator (2026)]{Advanced Validation of a High-Fidelity Hybrid Symplectic Propagator for Massive Space Debris Fragmentation Analysis}
\subtitle{Paper AAS 26-993}
\author[Puerta-Ibarra et al.]{Juan F. Puerta-Ibarra\inst{1} \and Bonnie Prado Pino\inst{2} \and Nicolas Mu\~noz-Galeano\inst{1}}
\institute[UdeA / Odyssey]{
  \inst{1}Universidad de Antioquia, Medell\'in, Colombia \and
  \inst{2}Odyssey Space Research LLC, Houston, TX, USA
}
\date{2026 AAS/AIAA Astrodynamics Specialist Conference}

\begin{document}

% ==============================================================================
% SLIDE 1: TITLE
% ==============================================================================
\begin{frame}
    \titlepage
    \begin{center}
        \footnotesize \textcolor{StructureGray}{Building upon computational methods established in AAS 25-787}
    \end{center}
\end{frame}

% ==============================================================================
% SLIDE 2: THE 2025 BASELINE & NEW NEED
% ==============================================================================
\begin{frame}{Introduction: Building Upon the 2025 Baseline}
    \begin{columns}[T]
        \begin{column}{0.55\textwidth}
            \begin{block}{AAS 25-787 Legacy Framework}
                \begin{itemize}
                    \item Validated conservative geopotential updates ($Central + J_2/J_4$) using high-performance WHFast routines.
                    \item Realized sub-kilometer coordinate agreement with NASA's GMAT across 20,000 unique LEO objects over a 24-hr cycle.
                \end{itemize}
            \end{block}
            \begin{alertblock}{The 2026 Dissipative Necessity}
                \begin{itemize}
                    \item Long-term (annual/decadal) tracking loops demand high-fidelity modeling of non-conservative forces.
                    \item Atmospheric drag and Solar Radiation Pressure (SRP) act as vital energy-dissipation constraints in LEO.
                \end{itemize}
            \end{alertblock}
        \end{column}
        \begin{column}{0.41\textwidth}
            \centering
            \vspace{0.2cm}
            \framebox[\linewidth]{%
                \begin{minipage}[c][4.2cm][c]{\linewidth}
                    \centering\footnotesize\textcolor{StructureGray}{
                        \includegraphics[width=\linewidth]{derivations/Space_debris_surrounding_Earth.jpg}
                    }
                \end{minipage}%
            }
        \end{column}
    \end{columns}
\end{frame}

% ==============================================================================
% SLIDE 3: SYMMETRIC STRANG SPLITTING
% ==============================================================================
\begin{frame}{Methodology: Second-Order Symmetric Strang Splitting}
    To incorporate non-Hamiltonian components without destabilizing symplectic qualitative conservation, the total evolution operator $\Phi(\Delta t)$ undergoes symmetric decomposition:
    \begin{equation*}
        \Phi(\Delta t) \approx \Phi_{P}\left(\frac{\Delta t}{2}\right) \circ \Phi_{G}(\Delta t) \circ \Phi_{P}\left(\frac{\Delta t}{2}\right)
    \end{equation*}
    
    \begin{columns}[T]
        \begin{column}{0.48\textwidth}
            \begin{block}{The Three-Phase Timestep ($\Delta t = 10$~s)}
                \begin{enumerate}
                    \item \textbf{Phase 1 ($\Phi_P$ Half-Step):} Velocity update via instantaneous drag and SRP accelerations.
                    \item \textbf{Phase 2 ($\Phi_G$ Full Step):} Stable integration of geopotential gravity via compiled WHFast core.
                    \item \textbf{Phase 3 ($\Phi_P$ Half-Step):} Final velocity adjustment using recomputed states.
                \end{enumerate}
            \end{block}
        \end{column}
        \begin{column}{0.48\textwidth}
            \centering
            \framebox[\linewidth]{%
                \begin{minipage}[c][3.8cm][c]{\linewidth}
                    \centering\footnotesize\textcolor{StructureGray}{
                        \includegraphics[width=\linewidth]{splitting_flowchart_v2.png}
                    }
                \end{minipage}%
            }
        \end{column}
    \end{columns}
\end{frame}

% ==============================================================================
% SLIDE 4: PHASE 1 — PINN THERMOSPHERIC SURROGATE
% ==============================================================================
\begin{frame}{Phase 1: Physics-Informed Neural Network (PINN) Surrogate}
    \begin{columns}[T]
        \begin{column}{0.52\textwidth}
            \begin{itemize}
                \item \textbf{The Lookup Bottleneck:} Piecewise exponential models create non-differentiable boundary steps ($\mathcal{C}^0$), injecting numerical noise.
                \item \textbf{The Solution:} Train an analytical network $\mathcal{N}_\theta: h \to \hat{\rho}$ regularized by the underlying physics PDE.
            \end{itemize}
            \begin{exampleblock}{Composite Loss Formulation}
                \[ \mathcal{L}_{Total}(\theta) = \mathcal{L}_{Data}(\theta) + \lambda \mathcal{L}_{Physics}(\theta) \]
                \[ \mathcal{L}_{Physics}(\theta) = \frac{1}{N_c} \sum_{j=1}^{N_c} \left\| \frac{d\mathcal{N}_\theta(h_j)}{dh} + \frac{\mathcal{N}_\theta(h_j)}{H(h_j)} \right\|^2 \]
            \end{exampleblock}
        \end{column}
        \begin{column}{0.44\textwidth}
            \centering
            \vspace{0.2cm}
            \framebox[\linewidth]{%
                \begin{minipage}[c][4.2cm][c]{\linewidth}
                    \centering\footnotesize\textcolor{StructureGray}{
                        \includegraphics[width=\linewidth]{derivations/PINN_Phase1.png}
                    }
                \end{minipage}%
            }
        \end{column}
    \end{columns}
\end{frame}

% ==============================================================================
% SLIDE 5: PHASE 1 — CONVERGENCE PROFILE
% ==============================================================================
\begin{frame}{Phase 1 Evaluation: Convergence \& Verification}
    \begin{columns}[c]
        \begin{column}{0.48\textwidth}
            \centering
            \textbf{Dual-Loss Log-Convergence Profile}
            \vspace{0.2cm}
            \framebox[\linewidth]{%
                \begin{minipage}[c][3.8cm][c]{\linewidth}
                    \centering\footnotesize\textcolor{StructureGray}{
                        \includegraphics[width=\linewidth]{derivations/fig1_pinn_loss.png}
                    }
                \end{minipage}%
            }
        \end{column}
        \begin{column}{0.48\textwidth}
            \centering
            \textbf{Vertical Density Field Validation}
            \vspace{0.2cm}
            \framebox[\linewidth]{%
                \begin{minipage}[c][3.8cm][c]{\linewidth}
                    \centering\footnotesize\textcolor{StructureGray}{
                        \includegraphics[width=\linewidth]{derivations/fig2_density_profile.png}
                    }
                \end{minipage}%
            }
        \end{column}
    \end{columns}
    \vspace{0.3cm}
    \centering
    \textcolor{UdeAGreen}{\textbf{Result:}} Joint optimization successfully guarantees continuous, smooth mass-decay derivatives without unphysical bounding ripples.
\end{frame}

% ==============================================================================
% SLIDE 6: PHASE 2 — GEOMETRIC PRUNING
% ==============================================================================
\begin{frame}{Phase 2: Breaking the Conjunction $\mathcal{O}(N^2)$ Trap}
    \begin{columns}[T]
        \begin{column}{0.55\textwidth}
            \begin{itemize}
                \item All-on-all baseline distance evaluations scale quadratically, introducing critical bottlenecks at $N=50,000$.
                \item \textbf{KD-Tree Domain Decomposition:} Partition the three-dimensional Euclidean state space recursively along alternating median planes.
            \end{itemize}
            \begin{block}{Localized Hyper-Spherical Filtering Query}
                \[ d(\mathbf{r}_{target}, \mathbf{r}_j) = \left\| \mathbf{r}_{target} - \mathbf{r}_j \right\|_2 \le R_{filter} \]
                \begin{itemize}
                    \item Isolates a tight proximal subset ($k \ll N$) within $R_{filter} = 50$~km constraints.
                    \item Truncates search complexity down to an asymptotically scalable $\mathcal{O}(N \log N)$ footprint.
                \end{itemize}
            \end{block}
        \end{column}
        \begin{column}{0.41\textwidth}
            \centering
            \vspace{0.4cm}
            \framebox[\linewidth]{%
                \begin{minipage}[c][3.8cm][c]{\linewidth}
                    \centering\footnotesize\textcolor{StructureGray}{
                        \includegraphics[width=\linewidth]{derivations/LSTM_Phase2.png}
                    }
                \end{minipage}%
            }
        \end{column}
    \end{columns}
\end{frame}

% ==============================================================================
% SLIDE 7: PHASE 2 — LSTM SEPARATION PERFORMANCE
% ==============================================================================
\begin{frame}{Phase 2 Evaluation: Conjunction Sequence Matrices}
    \begin{columns}[c]
        \begin{column}{0.48\textwidth}
            \centering
            \textbf{3D Spatial Clustered Neighborhood Selection}
            \vspace{0.2cm}
            \framebox[\linewidth]{%
                \begin{minipage}[c][3.8cm][c]{\linewidth}
                    \centering\footnotesize\textcolor{StructureGray}{
                        \includegraphics[width=\linewidth]{derivations/fig3_kdtree.png}
                    }
                \end{minipage}%
            }
        \end{column}
        \begin{column}{0.48\textwidth}
            \centering
            \textbf{LSTM Kinematic Arc Confusion Matrix}
            \vspace{0.2cm}
            \framebox[\linewidth]{%
                \begin{minipage}[c][3.8cm][c]{\linewidth}
                    \centering\footnotesize\textcolor{StructureGray}{
                        \includegraphics[width=\linewidth]{derivations/fig4_confusion_matrix.png}
                    }
                \end{minipage}%
            }
        \end{column}
    \end{columns}
    \vspace{0.3cm}
    \centering
    \textcolor{UdeAGreen}{\textbf{Result:}} The sequence architecture captures subtle variations in non-conservative energy decay with perfect recall despite massive class imbalances.
\end{frame}

% ==============================================================================
% SLIDE 8: PHASE 3 — INVERSE TARGETING VIA REINFORCEMENT LEARNING
% ==============================================================================
\begin{frame}{Phase 3: Event Reconstruction via Shared RAM Pointer Bridge}
    \begin{columns}[T]
        \begin{column}{0.54\textwidth}
            \begin{itemize}
                \item \textbf{The Inverse Problem:} Locate the initial coordinates, epoch ($\delta t$), and energy magnitude ($\delta v_{spread}$) of an uncataloged fragmentation center.
                \item \textbf{The Challenge:} Iterative file generation loops yield severe, disk-bound I/O serialization latencies.
                \item \textbf{The Shared Memory Solution:} Link the compiled C integrator directly into a Python Gymnasium MDP environment using \texttt{ctypes} pointers.
            \end{itemize}
            \begin{block}{Macro-Scale Moment Matching Reward}
                \[ R_t = - \frac{1}{C} \sqrt{ \left\| \boldsymbol{\mu}_{sim} - \boldsymbol{\mu}_{obs} \right\|_2^2 + \left\| \boldsymbol{\sigma}_{sim} - \boldsymbol{\sigma}_{obs} \right\|_2^2 } \]
            \end{block}
        \end{column}
        \begin{column}{0.42\textwidth}
            \centering
            \vspace{0.4cm}
            \framebox[\linewidth]{%
                \begin{minipage}[c][3.8cm][c]{\linewidth}
                    \centering\footnotesize\textcolor{StructureGray}{
                        \includegraphics[width=\linewidth]{derivations/RL_Phase3.png}
                    }
                \end{minipage}%
            }
        \end{column}
    \end{columns}
\end{frame}

% ==============================================================================
% SLIDE 9: PHASE 3 — PPO RESULTS BLOCK
% ==============================================================================
\begin{frame}{Phase 3 Evaluation: Pilot Gradient Trajectory Capture}
    \begin{columns}[c]
        \begin{column}{0.48\textwidth}
            \centering
            \textbf{PPO Initial Episode Loss Adaptation}
            \vspace{0.2cm}
            \framebox[\linewidth]{%
                \begin{minipage}[c][3.8cm][c]{\linewidth}
                    \centering\footnotesize\textcolor{StructureGray}{
                        \includegraphics[width=\linewidth]{derivations/fig5_rl_convergence.png}
                    }
                \end{minipage}%
            }
        \end{column}
        \begin{column}{0.48\textwidth}
            \centering
            \textbf{3D Pilot Reconstruction Tracking Vector}
            \vspace{0.2cm}
            \framebox[\linewidth]{%
                \begin{minipage}[c][3.8cm][c]{\linewidth}
                    \centering\footnotesize\textcolor{StructureGray}{
                        \includegraphics[width=\linewidth]{derivations/fig6_reconstruction.png}
                    }
                \end{minipage}%
            }
        \end{column}
    \end{columns}
    \vspace{0.3cm}
    \centering
    \textcolor{PhysicsOrange}{\textbf{Engineering Note:}} Truncated 100k pilot run acts as an architectural proof-of-concept. Scaling to a distributed HPC cluster is required to minimize final line residuals.
\end{frame}

% ==============================================================================
% SLIDE 10: ALGORITHMIC SCALABILITY CHART
% ==============================================================================
\begin{frame}{System Verification: Algorithmic Power-Law Scaling}
    \begin{columns}[T]
        \begin{column}{0.58\textwidth}
            \begin{block}{Preservation of Efficiency}
                \begin{itemize}
                    \item Complete pipeline benchmarked across 4 orders of magnitude ($N=50$ to $50,000$ independent fragments).
                    \item Avoids standard quadratic complexity explosions ($\mathcal{O}(N^2)$), tracking strictly parallel to the geopotential baseline.
                    \item Confirms a highly scalable power-law performance exponent of approximately $\mathcal{O}(N^{1.60})$.
                \end{itemize}
            \end{block}
            \centering
            \vspace{0.1cm}
            \textcolor{UdeAGreen}{\textbf{Core Mechanism:}} Parallel vectorized tensor evaluations coupled with localized tree domain pruning successfully absorb non-conservative physical latencies.
        \end{column}
        \begin{column}{0.38\textwidth}
            \centering
            \vspace{0.2cm}
            \framebox[\linewidth]{%
                \begin{minipage}[c][4.0cm][c]{\linewidth}
                    \centering\footnotesize\textcolor{StructureGray}{
                        \includegraphics[width=\linewidth]{derivations/PowerLaw2026.png}
                    }
                \end{minipage}%
            }
        \end{column}
    \end{columns}
\end{frame}

% ==============================================================================
% SLIDE 11: CRITICAL ARCHITECTURAL TRADE-OFF MATRIX
% ==============================================================================
\begin{frame}{Architectural Trade-Off Analysis: Framework Pros \& Cons}
    \begin{table}
        \small
        \centering
        \begin{tabular}{p{0.46\textwidth} p{0.48\textwidth}}
            \toprule
            \textbf{Engineering Strengths (Pros)} & \textbf{Structural Constraints (Cons)} \\
            \midrule
            Maintenance of the baseline $\mathcal{O}(N^{1.60})$ power-law processing slope. & Constant vertical latency shift due to continuous network model inference overhead. \\ \addlinespace
            Continuous, $\mathcal{C}^\infty$-differentiable fields completely bypass interpolation discontinuities. & Crossover threshold ($N \approx 500$) below which standard analytical integration is superior. \\ \addlinespace
            Direct pointer memory bridge eliminates slow file-serialization bottlenecks. & Massive array allocation within RAM introduces workstation memory boundaries at $N \ge 500,000$. \\ \addlinespace
            Macro-scale moment calculations prune combinatorial association demands. & Out-of-domain performance degradation when subjected to uncalibrated solar storm fluxes. \\
            \bottomrule
        \end{tabular}
    \end{table}
\end{frame}

% ==============================================================================
% SLIDE 12: CONCLUSIONS & ROADMAP
% ==============================================================================
\begin{frame}{Conclusions \& Operational Horizons}
    \begin{columns}[T]
        \begin{column}{0.48\textwidth}
            \begin{block}{Core Achievements}
                \begin{itemize}
                    \item Reconciled semi-symplectic numerical integration stability with dissipative physics using second-order Strang splitting.
                    \item Maintained strict physical limits on relative secular energy drifts ($10^{-8}$ to $10^{-7}$).
                    \item Proved the feasibility of localized neural surrogates and memory-cached pointer arrays.
                \end{itemize}
            \end{block}
        \end{column}
        \begin{column}{0.48\textwidth}
            \begin{exampleblock}{The Next Frontiers}
                \begin{itemize}
                    \item \textbf{HPC Cluster Scaling:} Transition from local workstation cores to distributed multi-node clusters to reach $10^7$ optimization horizons.
                    \item \textbf{Attention-Based Models:} Integrate spatial-temporal Transformers to handle multi-body cislunar gravity fields.
                    \item \textbf{Temporal Inversion:} Exploit time-reversible flow parameters for true autonomous backward-in-time orbital propagation.
                \end{itemize}
            \end{exampleblock}
        \end{column}
    \end{columns}
    \vspace{0.4cm}
    \centering
    \footnotesize \textbf{Contact:} \texttt{juanf.puerta@udea.edu.co} \, | \, \textbf{Repository:} \texttt{[INSERT: GitHub / Presentation Link QR Box]}
\end{frame}

\end{document}